\scanpage{143}
% archive foliation: 142
\deleted{Ceci prouve}

D'autre part on a
\[
\unsure{\mathfrak{Z}}(G\times G')=
\unsure{\mathfrak{Z}}(G)\times\unsure{\mathfrak{Z}}(G')
\]
et l'ens. des sousgroupes $\in C$ d'un $G\in C$ satisfait la condition minimale.

Ces propriétés formelles suffisent pour construire un \uline{proobjet} de $C$,
au moyen d'un système projectif de épimorphismes, qu'on
\deleted{\unsure{va}} désignera par $\pi_1^C(S,\xi)$.

\uline{Exemple I} $C=$ \unsure{ens|cat} des groupes définis par des
groupes finis ordinaires. On retrouve les groupes fondamentaux.

\uline{Exemple II} $S$ est un préschéma au dessus d'un corps $k$.
$C$ provient d'une catégorie $C_0$ de $k$-groupes finis
(par exple \deleted{\uline{\unsure{tous|\illegible{4}}}} les $k$-groupes finis --- ce qui
\deleted{\unsure{est}}\inserted{semble}
raisonnable surtout si $k$ sep. clos)
stable par produits, noyaux de couples, images.
